
\documentclass[../main]{subfiles}
\usepackage{graphicx} % Para incluir imágenes

\begin{document}

\chapter{Ciclo de Carnot}

\section{Introducción}

El ciclo de Carnot es un ciclo termodinámico ideal que describe el funcionamiento de una máquina térmica reversible. Es el ciclo más eficiente que puede operar entre dos reservorios de calor a diferentes temperaturas. El ciclo de Carnot tiene una gran importancia teórica y práctica, ya que establece el límite de eficiencia para cualquier máquina térmica.

\section{Teoría}

El ciclo de Carnot consta de cuatro procesos reversibles:

1. **Expansión isotérmica:** El gas absorbe calor del reservorio caliente a temperatura constante $T_H$. El volumen del gas aumenta y su presión disminuye.

2. **Expansión adiabática:** El gas se expande sin intercambio de calor con el entorno. La temperatura del gas disminuye.

3. **Compresión isotérmica:** El gas cede calor al reservorio frío a temperatura constante $T_C$. El volumen del gas disminuye y su presión aumenta.

4. **Compresión adiabática:** El gas se comprime sin intercambio de calor con el entorno. La temperatura del gas aumenta.

La eficiencia del ciclo de Carnot, representada por la letra griega $\eta$, se define como la relación entre el trabajo realizado por la máquina y el calor absorbido del reservorio caliente:

$$\eta = \frac{W}{Q_H} = 1 - \frac{T_C}{T_H}$$

Donde:

* $W$ es el trabajo realizado por la máquina
* $Q_H$ es el calor absorbido del reservorio caliente
* $T_C$ es la temperatura del reservorio frío
* $T_H$ es la temperatura del reservorio caliente

**Ejemplo:**

Consideremos un ciclo de Carnot que opera entre un reservorio caliente a 500 K y un reservorio frío a 300 K. La máquina absorbe 1000 J de calor del reservorio caliente. ¿Cuál es el trabajo realizado por la máquina?

**Solución:**

Primero, calculamos la eficiencia del ciclo:

$$\eta = 1 - \frac{T_C}{T_H} = 1 - \frac{300 K}{500 K} = 0.4$$

Luego, calculamos el trabajo realizado por la máquina:

$$W = \eta Q_H = 0.4 \times 1000 J = 400 J$$

**Fórmulas importantes:**

* Eficiencia del ciclo de Carnot: $\eta = 1 - \frac{T_C}{T_H}$
* Calor absorbido del reservorio caliente: $Q_H = \frac{W}{\eta}$
* Calor cedido al reservorio frío: $Q_C = Q_H - W$

\section{Cuestionario}

1. ¿Qué es el ciclo de Carnot?

2. ¿Cuáles son los cuatro procesos reversibles que componen el ciclo de Carnot?

3. ¿Cómo se define la eficiencia del ciclo de Carnot?

4. ¿De qué depende la eficiencia del ciclo de Carnot?

5. ¿Cuál es el límite de eficiencia para cualquier máquina térmica?

6. ¿En qué tipo de máquinas se utiliza el ciclo de Carnot?

7. ¿Cuáles son las ventajas y desventajas del ciclo de Carnot?

8. ¿Qué es la entropía y cómo se relaciona con el ciclo de Carnot?

**Nota:** El código LaTeX se ha dividido en dos partes debido a que supera los 4096 caracteres. La segunda parte se encuentra en el siguiente comentario.

\end{document}